\documentclass[12pt,a4paper]{article}
\usepackage[utf8]{inputenc}
\usepackage[vietnamese]{babel}
\usepackage[left=2cm,right=2cm,top=2cm,bottom=2cm]{geometry}
\usepackage{mathptmx}
\usepackage{enumitem}

\pagestyle{plain}
\linespread{1.15}

\begin{document}

% --- HEADER BÁO CHÍ LƯU TRỮ ---
\noindent
\begin{minipage}[t]{0.45\textwidth}
    \begin{center}
        \textbf{\large BÁO CÔNG AN THỦ ĐÔ}\\
        \textit{Cơ quan ngôn luận Công an TP. Hà Nội}\\[0.05cm]
        \small Số ra ngày: 14 tháng 07 năm 1998
    \end{center}
\end{minipage}
\hfill
\begin{minipage}[t]{0.52\textwidth}
    \begin{center}
        \small
        \textbf{CHUYÊN MỤC TRẬT TỰ \& ĐỜI SỐNG}\\
        \textit{Trang 04 — Tin tức địa phương}\\[0.05cm]
        Mã hồ sơ lưu trữ: \texttt{CATD-1998-T07-042}
    \end{center}
\end{minipage}

\vspace{0.4cm}
\hrule
\vspace{0.4cm}

% --- TIÊU ĐỀ BÀI BÁO ---
\begin{center}
    \textbf{\Large SỰ CỐ THƯƠNG TÂM: BÉ TRAI TỬ VONG DO NGẠT KHÍ KHI BỊ KẸT TRONG TỦ GỖ CŨ}\\
    \textit{(Vụ tai nạn đau lòng xảy ra tại khu đất trống Phân khu Cảng)}
\end{center}

\vspace{0.2cm}

\noindent \textbf{HÀ NỘI —} Chiều ngày 12/07/1998, tại khu vực bãi đất trống giáp ranh tổ 4 Phân khu Cảng, người dân địa phương bàng hoàng phát hiện thi thể cháu bé **Nguyễn Gia Huy (sinh năm 1990 — 8 tuổi)** tử vong bên trong một chiếc tủ gỗ ép thanh lý bỏ hoang.

\section*{DIỄN BIẾN SỰ VIỆC BAN ĐẦU}
Theo điều tra ban đầu của Công an cơ sở, vào sáng cùng ngày, nhóm trẻ em trong xóm rủ nhau chơi trò trốn tìm quanh khu vực bãi tập kết phế liệu. Cháu Huy đã chui vào ẩn nấp trong chiếc tủ gỗ cũ. 

Đáng tiếc, cánh cửa tủ có cấu tạo chốt cài tự sập từ phía ngoài. Khi cánh cửa đóng sập lại, chốt gỗ bên ngoài rơi xuống khóa chặt khiến cháu bé không thể tự đẩy cửa thoát ra ngoài. Được biết, **cháu Huy vốn bị câm bẩm sinh và có tiền sử bệnh tim**, do quá hoảng sợ trong không gian tối kín khí, cháu đã lên cơn khó thở cấp dẫn đến ngạt khí và tử vong thương tâm trước khi gia đình tìm thấy.

\section*{LỜI KỂ CỦA NHỮNG ĐỨA TRẺ CÙNG XÓM}
Phóng viên ghi nhận chia sẻ từ những đứa trẻ cùng tham gia buổi chơi hôm đó:
\begin{itemize}[leftmargin=1.5cm]
    \item \textbf{Cháu M. (Trần Ngọc Mai — 6 tuổi):} Hốt hoảng bật khóc: \textit{"Cháu không hiểu sao bạn Huy lại tự trốn vào tủ sâu thế... Bạn Huy bình thường nhút nhát và chậm chạp lắm, lúc nào chơi bọn cháu cũng phải dắt tay chỉ chỗ..."}
    \item \textbf{Cháu K. (Nguyễn Văn Khang — 8 tuổi):} Đứng trước ống kính phóng viên, cháu K. tỏ vẻ mặt buồn bã: \textit{"Lúc chơi cháu đi tìm mãi không thấy bạn Huy đâu... Cháu không ngờ bạn ấy lại chui vào chiếc tủ đó. Cháu thương bạn ấy lắm..."}
\end{itemize}

\section*{LỜI CẢNH TỈNH CHO CÁC BẬC PHỤ HUYNH}
Cơ quan chức năng khuyến cáo các gia đình có con nhỏ cần hết sức lưu ý trông nom, không để trẻ em tự do vui chơi tại các khu vực bãi hoang, công trình phế liệu nguy hiểm.

\vspace{0.6cm}
\noindent \textit{(Ghi chú nghiệp vụ 2026: Đây là toàn văn nội dung bài báo sau khi Đội khám nghiệm thu thập và ghép hoàn chỉnh các mảnh báo bị xé vụn ký hiệu \texttt{p5} rơi tại hiện trường nhà Khang đêm 24/07/2026).}

\end{document}
